\documentclass{beamer}

\usetheme{Madrid} % Tu peux changer : Berlin, CambridgeUS, Warsaw, etc.

\title{Image Restoration by Inverse Problem with Kernel Estimation}
\author{Rothlingshofer Yanic \\ Douar Lounes}
\date{21 novembre 2025}

\begin{document}

% Page de titre
\begin{frame}
    \titlepage
\end{frame}

% Plan
\begin{frame}{Plan}
    \tableofcontents
\end{frame}

\section{Introduction}

\begin{frame}{Context}
    \begin{itemize}
        \item Motion blur occurs when camera or scene motion affects the exposure.
        \item Variational and Bayesian approaches historically dominate blind deblurring.
        \item Deep learning achieves strong results but requires large datasets and heavy computation.
        \item We focus on an interpretable, training-free variational pipeline based on spectral priors.
    \end{itemize}
\end{frame}

\begin{frame}{Objectives}
    \begin{itemize}
        \item Estimate the blur kernel from the blurred image using anisotropic power spectrum modelling.
        \item Recover the sharp image through TV-regularised deconvolution solved by Split-Bregman.
        \item Pipeline:
        \begin{itemize}
            \item spectral modulus estimation of the kernel,
            \item phase reconstruction under support and positivity constraints,
            \item TV-based restoration with tuned hyperparameters.
        \end{itemize}
    \end{itemize}
\end{frame}

\section{Mathematical Model}

\begin{frame}{Blur Model}
    On considère :
    \[
    v = u * k + b
    \]
    \begin{itemize}
        \item $u$: sharp latent image.
        \item $k$: non-negative, normalized blur kernel with compact support.
        \item $b$: additive impulse noise.
        \item Ill-posed inverse problem requiring prior assumptions.
        \item We assume spatially invariant linear motion blur.
    \end{itemize}
\end{frame}


\section{Blind Kernel Estimation}

\begin{frame}{Prior}
\[
| \widehat{u}(\xi n_\theta) |^2 \approx \frac{1}{\varepsilon + c_\theta \|\xi\|^2}
\]
\[
\Rightarrow | \widehat{D_\theta u}(\xi n_\theta) |^2 = |\xi|^2 \, | \widehat{u}(\xi n_\theta)|^2 = \frac{|\xi|^2}{\varepsilon + c_\theta \|\xi\|^2}
\]
\begin{itemize}
    \item Natural images follow an anisotropic power-law decay.
    \item Taking directional derivatives amplifies high frequencies.
    \item This allows isolating the blur kernel in the spectrum of $D_\theta v$.
\end{itemize}
\end{frame}

\begin{frame}{Whitening and Directional Derivatives}
    \[
    D_\theta v = (D_\theta u) * k + D_\theta b
    \]
    \[
    \widehat{D_\theta v}(\xi n_\theta) = i(\xi n_\theta \cdot n_\theta)\, \widehat{u}(\xi n_\theta)\, \widehat{k}(\xi n_\theta) + \widehat{D_\theta b}(\xi n_\theta)
    \]
    \begin{itemize}
        \item Whitening: $D_\theta$ flattens the spectrum of natural images.
        \item The remaining directional variations come from the blur kernel.
        \item Enables isolating $|\widehat{k}|^2$ in the autocorrelation domain.
    \end{itemize}

\end{frame}

\begin{frame}{Autocorrelation Analysis}
    \[
    \mathbb{E}\left(|\widehat{D_\theta v}(\xi n_\theta)|^2\right) 
    = |\xi|^2 \, |\widehat{u}(\xi n_\theta)|^2 \, |\widehat{k}(\xi n_\theta)|^2 
    + \mathbb{E}\left(|\widehat{D_\theta b}|^2\right)
    \]
    \[
    = c_\theta |\widehat{k}(\xi n_\theta)|^2 + \mu_\theta
    \]
    \begin{itemize}
        \item After whitening, the residual spectrum becomes linear in $|\widehat{k}|^2$.
        \item $\mu_\theta$ is estimated using the kernel normalization constraint:
        \[
        \int_\mathbb{R} P_\theta(h)(x)\,dx = 1.
        \]
        \item Allows reconstruction of $|\widehat{k}(\xi)|^2$ for all $\theta$.
    \end{itemize}
\end{frame}


\begin{frame}{Autocorrelation Deconvolution with compensation filter}
    \[
    R_{D_{\theta} v} = R_{D_{\theta} u} * R_k + R_{D_{\theta} b}
    \]
    \begin{itemize}
        \item Autocorrelation of $D_\theta v$ gives access to the width of the blur in direction $\theta$.
        \item A compensation filter deconvolves camera noise and intrinsic blur using CG iterations.
        \item Negative lobes near the center invalidate the compensation and revert to the raw autocorrelation.
    \end{itemize}
\end{frame}

\begin{frame}{Initial Support Estimation}
    \begin{itemize}
        \item Initial support $s_\theta$ estimated from minima of compensated autocorrelation.
        \item Support is smoothed using the Lipschitz constraint:
        \[
        |s_{\theta_i} - s_{\theta_j}| \le \kappa |\theta_i - \theta_j|.
        \]
        \item Ensures coherent support across all projection directions.
    \end{itemize}
\end{frame}

\begin{frame}{Power spectrum estimation}
    \[
    F(R(D_\theta v))(\xi) = c_\theta |\widehat{k}(\xi n_\theta)|^2 + \mu_\theta
    \]
    \[
    \Rightarrow |\widehat{k}(\xi n_\theta)|^2 = \frac{F(R(D_\theta v))(\xi) - \mu_\theta}{c_\theta}
    \]
    \begin{itemize}
        \item Full modulus $|\widehat{k}(\xi)|$ recovered by interpolating directional slices.
        \item Noise bias removed using low-frequency estimates of $\mu_\theta$.
        \item Produces a stable modulus usable for phase retrieval.
    \end{itemize}
\end{frame} 

\begin{frame}{Phase retrieval and kernel estimation}
    \begin{itemize}
        \item Random phases $\phi(\xi)$ sampled to reconstruct candidate kernels:
        \[
        k = \mathcal{F}^{-1}\left(|\widehat{k}(\xi)| e^{i\phi(\xi)}\right)
        \]
        \item Constraints applied: positivity, compact support, normalization.
        \item Best candidate chosen by TV-deconvolution PSNR score.
    \end{itemize}
\end{frame}


\begin{frame}{TV Deconvolution}
    \[
    \hat{u} = \arg\min_u 
    \frac12 \|k*u - v\|_2^2 
    + \lambda \|\nabla u\|_1
    \]
    \[
    \text{Split-Bregman iterations:}
    \]
    \[
    u^{k+1} = \arg\min_u 
    \frac12\|k*u - v\|_2^2 
    + \frac{\gamma}{2}\|\nabla u - d^k + b^k\|_2^2
    \]
    \[
    d^{k+1} = \operatorname{shrink}(\nabla u^{k+1} + b^k,\, \lambda/\gamma)
    \]
    \[
    b^{k+1} = b^k + (\nabla u^{k+1} - d^{k+1})
    \]
    \begin{itemize}
        \item Solved in Fourier domain using symmetric or circular boundary conditions.
        \item TV prior preserves edges and limits ringing.
        \item Convergence when $\|u^{k+1}-u^k\| < \eta$.
    \end{itemize}
\end{frame}

\section{Kernel Estimation}

\begin{frame}{Estimation Principle}
    \begin{itemize}
        \item Analysis of the directional spectrum of the image.
        \item Use of directional derivatives.
        \item Recovery of the Fourier transform modulus of the kernel.
        \item Kernel reconstruction via random phases.
    \end{itemize}
\end{frame}

\section{Results}

\begin{frame}{Types of Experiments Conducted}
    \begin{itemize}
        \item Validation with known vs estimated blur.
        \item Robustness to noise.
        \item Influence of boundary conditions.
        \item Effect of kernel size.
        \item Tests on real images.
    \end{itemize}
\end{frame}

\end{document}